\section{总结与延伸}

这堂课最重要的贡献不是给 DPO、PPO 或 QLoRA 排名，而是提供一段可解释的开放对齐系统史：每次模型跃迁都同时依赖 base model、数据分布、训练可访问性、评测反馈和发布 artifact。理解这条链，才能判断一个新方法究竟解决了哪个瓶颈。

\subsection{十四条核心结论}

本节把课程压缩成可复用结论，每条都保留条件和证据等级。

\begin{enumerate}
  \item RLHF 对 ChatGPT-like behavior 很重要，但不是完整产品的充分条件。
  \item Base model 是能力底座，alignment 主要重新排序和塑造交互分布。
  \item IFT、SFT、RLHF 与 DPO 必须按数据和损失区分。
  \item Alpaca 证明 synthetic instruction data 可快速适配接口。
  \item Vicuna/ShareGPT 证明真实用户 prompt distribution 极有价值。
  \item OpenAssistant 说明开放人类数据具有长期复用价值。
  \item QLoRA 扩大参与者范围，但没有自动解决 RLHF 稳定性。
  \item Safety 与 uncensored 争论本质上是拒答边界和用户控制的选择。
  \item Arena、AlpacaEval、MT-Bench 与 Leaderboard 服务不同反馈时长。
  \item LLM-as-a-judge 快但有 length、position、self-preference 和 contamination bias。
  \item RLHF 目标结合 reward 与 KL constraint；PPO 只是实现路线之一。
  \item DPO 简单且有效，但与 PPO 的训练动态和数据利用不同。
  \item Zephyr/Tulu 证明 DPO 可落地，SteerLM/Starling 证明 RL 仍有竞争力。
  \item 2024 年开放 alignment 的稳定瓶颈是多样、合法、可发布的 preference data。
\end{enumerate}

\subsection{一张表串起开放对齐系统}

为了避免只记模型名，下面按系统组件总结每一阶段解决的问题与残余风险。

\begin{table}[H]
\centering
\small
\begin{tabularx}{\textwidth}{L{0.16\textwidth} L{0.22\textwidth} X}
\toprule
组件 & 代表 artifact & 仍未解决的问题 \\
\midrule
Base model & LLaMA/Llama 2/Llama 3 & 能力、许可、数据 provenance 与训练成本 \\
IFT data & Alpaca、ShareGPT、UltraChat & 真实 intent、多样性、隐私与模板偏差 \\
Human data & OpenAssistant/Arena logs & 成本、一致性、consent 与覆盖 \\
Efficient tuning & LoRA/QLoRA & 完整能力保留、RL stability 与 inference 合并 \\
Evaluation & Arena/AlpacaEval/MT-Bench/Leaderboard & judge bias、污染、延迟与 gaming \\
Preference optimization & PPO/DPO/SteerLM & 数据效率、稳定性、reward hacking 与可解释性 \\
Ecosystem & AI2、HF、LMSYS、community & lineage、可复现性、长期维护与治理 \\
\bottomrule
\end{tabularx}
\caption{开放对齐的组件、代表 artifact 与残余风险。}
\end{table}

\subsection{实践用 release checklist}

课程历史可以转成一个模型发布前 checklist。目标是让“我们用了 DPO”不再替代完整证据。

\begin{lstlisting}[caption={Open aligned model release checklist}]
record_base_model_and_license()
publish_chat_template_and_training_objective()
document_instruction_and_preference_data_provenance()
report_compute_memory_and_random_seeds()
run_capability_safety_and_human_preference_evals()
disclose_judge_models_baselines_and_known_biases()
release_failure_examples_and_model_card()
\end{lstlisting}

\begin{importantbox}{判断新 alignment 方法的五道门}
\begin{enumerate}
  \item 是否在共享 base/data/eval 下超过强 baseline？
  \item 是否产生可下载、可复现的竞争模型，而不只是训练曲线？
  \item 是否保持 base capabilities，并降低真实 failure rate？
  \item 是否在多个 judge 和人类偏好上成立，而非单一 leaderboard？
  \item 收益是否值得额外数据、rollout、显存与工程复杂度？
\end{enumerate}
\end{importantbox}

\subsection{自测问题}

以下问题用于检查是否掌握系统因果，而不是只记 release 时间线。

\begin{enumerate}
  \item 为什么 RLHF 对 ChatGPT 重要，却仍不是充分条件？
  \item IFT 与 SFT 在目标、数据和接口上如何区分？
  \item ShareGPT 为什么可能比更多 synthetic prompts 更有价值？
  \item Weight differences 与 LoRA update 有什么本质差异？
  \item QLoRA 节省了哪些内存，又没有节省哪些 RLHF 成本？
  \item Arena、AlpacaEval、MT-Bench 和 Leaderboard 分别适合什么决策？
  \item KL penalty 在 RLHF objective 中防止什么？
  \item DPO loss 中 policy/reference log-ratio 表示什么？
  \item 为什么 Zephyr/Tulu 不能单独证明 DPO 总是优于 PPO？
  \item 合成 preference data 为什么可能扩大规模却缩小分布？
\end{enumerate}

\subsection{拓展阅读}

\begin{itemize}
  \item Ouyang et al., \emph{Training language models to follow instructions with human feedback}：InstructGPT 与 RLHF pipeline。
  \item Taori et al., \emph{Stanford Alpaca}；Wang et al., \emph{Self-Instruct}：早期 synthetic instruction data。
  \item Dettmers et al., \emph{QLoRA}：量化 base model 与低秩 fine-tuning。
  \item Rafailov et al., \emph{Direct Preference Optimization}：DPO 推导与实验。
  \item Bai et al., \emph{Constitutional AI}；Köpf et al., \emph{OpenAssistant Conversations}：偏好与开放人类数据。
  \item Zheng et al., \emph{Chatbot Arena}：人类 pairwise evaluation 与开放模型生态。
\end{itemize}

\end{document}
